%
\hsection{While-Loop / While-Schleife}%
%
%
\begin{question}{What is it? / Was ist das?}%
\begin{itemize}%
\item[EN] Explain the syntax \emph{and} purpose/use of an \pythonilIdx{while} statement in \python. %
You can use an example for your explanation.%
\item[DE] Erklären Sie die Syntax \emph{und} den Verwendungszweck eines \pythonil{while}-Statements in \python. %
Sie können ein Beispiel in Ihrer Erklärung verwenden.%
\end{itemize}%
\end{question}%
%
%
\begin{question}{Program Understanding / Programm Verstehen}%
\gitLoadPython{while_01}{}{questions/while_01.py}{}%
\listingPython{while_01}{The program \programUrl{while_01}}%
\begin{itemize}%
\item[EN] Explain program \programUrl{while_01} in \cref{lst:while_01} \emph{line-by-line}, i.e., every single line. %
What output does it produce?
\item[DE] Erklären Sie das Programm \programUrl{while_01} in \cref{lst:while_01} Zeile für Zeile, also jede einzelne Zeile. %
Welche Ausgabe produziert es?%
\end{itemize}%
\end{question}%
%
%
%
\begin{question}{Programming / Programmieren}%
\begin{itemize}%
\item[EN] The variable \pythonil{text} holds text string. %
Write some \python\ code that iterates over the characters in \pythonil{text} until it either reaches the character~\pythonil{\"Q\"} or the end of the string and that sums up all digits and prints the result. %
For example, if \pythonil{text = \"ksd934.df4Qff4\"}, it would print~\textil{20} and for \pythonil{text = \"f5q6GH\"}, it would print~\textil{11}.%
%
\item[DE] Die Variable \pythonil{text} speichert einen Text-String. %
Schreiben Sie \python-Kode der über die Zeichen in \pythonil{text} iteriert bis er entweder das Zeichen~\pythonil{\"Q\"} oder das Ende des Strings erreicht und alle Ziffern, die er findet, aufaddiert und das Ergebnis ausgibt. %
Zum Beispiel für \pythonil{text = \"ksd934.df4Qff4\"} würde er~\textil{20} ausgeben und für \pythonil{text = \"f5q6GH\"} würde er~\textil{11} ausgeben.%
\end{itemize}%
\end{question}%
%
%
\begin{question}{Programming / Programmieren}%
\begin{equation}%
a_n=\left\{%
\begin{array}{lll}%
\frac{1}{2} a_{n-1}&\text{if }a_{n-1}\text{ even}&\text{wenn }a_{n-1}\text{ gerade}\\%
3 a_{n-1} + 1&\text{if }a_{n-1}\text{ odd}&\text{wenn }a_{n-1}\text{ ungerade}%
\end{array}%
\right.%
\end{equation}%
\begin{itemize}%
\item[EN] The sequence defined by the equation above is called Collatz Problem~\cite{L2010TUCT3P}. %
It is believed that, for any positive integer $a_0\in\naturalNumbersO$ with~$a_0>0$, there will eventually be an $a_n=1$. %
Expect that $a_0$ be an integer greater than zero stored in a variable~\pythonil{a}. %
Write \python\ code that prints all the values of the corresponding Collatz Problem sequence until (and including)~\pythonil{1}.%
%
\item[DE] Die Sequence die die Gleichung oben definiert wird Collatz Problem genannt~\cite{L2010TUCT3P}. %
Es wird angenommen, dass für jede positive Ganzzahl $a_0\in\naturalNumbersO$ mit~$a_0>0$ irgendwann ein $a_n=1$ erreicht wird. %
Nehmen Sie an, dass $a_0$ eine Ganzzahl größer als Null gespeichert in der Variable~\pythonil{a} ist. %
Schreiben Sie \python-Kode der alle Werte der entsprechenden Collatz-Problem-Sequence bis einschließlich~\pythonil{1} ausgibt.%
\end{itemize}%
\end{question}%
%
%
\begin{question}{Programming / Programmieren}%
\begin{itemize}%
\item[EN] Your bank account at the \emph{Bank of Python} holds \pythonil{money = 20_000}~RMB. %
Every year, you get~\pythonil{interest}\% of interest, i.e., after one year, you have $\left\lfloor(1+\pythonil{interest}/100)\pythonil{money}\right\rfloor$~RMB.
Notice that $\lfloor M\rfloor$ means rounding down to the nearest integer smaller or equal to~$M$.
Write \python\ code that counts the number of years you have to wait until there are $100\decSep0000$~(=~100万)~RMB in your account~(if you do not deposit or withdraw any money).%
%
\item[DE] Ihr Bankkonto bei der \emph{Bank of Python} hat \pythonil{money = 20_000}~RMB. %
Jedes Jahr beommen Sie \pythonil{interest}\%~Zinsen.
Also nach einem Jahr haben Sie $\left\lfloor(1+\pythonil{interest}/100)\pythonil{money}\right\rfloor$~RMB.
Beachten Sie dass $\lfloor M\rfloor$ bedeutet, das zur nächsten Ganzzahl kleiner oder gleich~$M$ (ab)gerundet wird.
Schreiben Sie \python-Kode der die Anzahl von Jahren zählt, die Sie warten müssen, bis $100\decSep0000$~(=~100万)~RMB in Ihrem Konto sind~(wenn Sie weder Geld einzahlen noch abheben).%
\end{itemize}%
\end{question}%
%
\endhsection%
%
